\section{BUG-023: Audit Log File Grows Indefinitely Without Rotation}
\label{sec:bug-023}

\subsection*{Metadata}
\begin{tabular}{ll}
\textbf{Issue ID} & BUG-023 \\
\textbf{Severity} & Medium \\
\textbf{Category} & Bug Fix \\
\textbf{Affected File} & \srcfile{src/audit.ts:37-42} \\
\textbf{Branch} & \branchref{fix/BUG-023-audit-log-rotation} \\
\textbf{Changes} & \branchdiff{fix/BUG-023-audit-log-rotation} \\
\end{tabular}

\subsection*{Executive Summary}
The \texttt{AuditLogger} class in \texttt{src/audit.ts:16-68} appends every event as a JSON line to a single file (\texttt{audit.jsonl}) without any size limit, rotation, or truncation. For long-running agent sessions — particularly the voice server (\texttt{src/voice/server.ts}), which can run for hours or days — this file grows without bound, consuming all available disk space.

Each tool call generates at minimum two audit entries (one for \texttt{logToolUse}, one for \texttt{logToolResult}), plus \texttt{logResponse} after each turn. For an active agent session making hundreds or thousands of tool calls, the audit file can easily grow to hundreds of megabytes. On systems with limited disk space (cloud containers, Raspberry Pi, ephemeral CI environments), this can cause the process to crash with \texttt{ENOSPC} (no space left on device) errors, or silently corrupt the audit file itself if a write is truncated mid-line.

The root cause is twofold: (1) \texttt{appendFile} is called without any size check before writing, and (2) there is no rotation mechanism — neither time-based (daily) nor size-based (per-N-bytes). The class has no configuration options for max file size, retention count, or compression.

---

\subsection*{Root Cause Analysis}
\#\#\# The Code



\begin{lstlisting}[language=TypeScript]
// audit.ts:16-68 — AuditLogger class
export class AuditLogger {
	private logPath: string;
	private sessionId: string;
	private enabled: boolean;

	constructor(gitagentDir: string, sessionId: string, enabled: boolean) {
		this.logPath = join(gitagentDir, "audit.jsonl");
		this.sessionId = sessionId;
		this.enabled = enabled;
	}

	async log(event: string, data: Partial<AuditEntry> = {}): Promise<void> {
		if (!this.enabled) return;

		const entry: AuditEntry = {
			timestamp: new Date().toISOString(),
			session_id: this.sessionId,
			event,
			...data,
		};

		try {
			await mkdir(dirname(this.logPath), { recursive: true });
			await appendFile(this.logPath, JSON.stringify(entry) + "\n", "utf-8");  // ← Always appends, never checks size
		} catch {
			// Audit logging failures are non-fatal
		}
	}

	async logToolUse(tool: string, args: Record<string, any>): Promise<void> {
		await this.log("tool_use", { tool, args });
	}

	async logToolResult(tool: string, result: string): Promise<void> {
		await this.log("tool_result", { tool, result: result.slice(0, 1000) });
	}

	async logResponse(): Promise<void> {
		await this.log("response");
	}

	async logError(error: string): Promise<void> {
		await this.log("error", { error });
	}

	async logSessionStart(): Promise<void> {
		await this.log("session_start");
	}

	async logSessionEnd(): Promise<void> {
		await this.log("session_end");
	}
}

\end{lstlisting}



\#\#\# Growth Projection

Each audit entry is approximately 200-500 bytes. For an active session:

| Metric | Per Turn | Per 100 Turns | Per 1000 Turns |
|---|---|---|---|
| \texttt{logToolUse} (1/tool) | \textasciitilde{}300 bytes | 30 KB | 300 KB |
| \texttt{logToolResult} (1/tool) | \textasciitilde{}1.2 KB (result truncated to 1000 chars) | 120 KB | 1.2 MB |
| \texttt{logResponse} (1/turn) | \textasciitilde{}200 bytes | 20 KB | 200 KB |
| \textbf{Total per turn} | \textasciitilde{}1.7 KB | 170 KB | 1.7 MB |
| \textbf{Total for 10k turns} | 17 MB | — | — |
| \textbf{Total for 100k turns} | 170 MB | — | — |
| \textbf{Total for 1M turns} | 1.7 GB | — | — |

A voice server session handling 100k turns (not unrealistic for a 24-hour deployment) generates \textasciitilde{}170 MB of audit data. In a container with a 1 GB ephemeral disk, this can fill the disk in \textasciitilde{}6 hours.

\#\#\# The \texttt{catch} Block Swallows Disk-Full Errors



\begin{lstlisting}[language=TypeScript]
try {
	await mkdir(dirname(this.logPath), { recursive: true });
	await appendFile(this.logPath, JSON.stringify(entry) + "\n", "utf-8");
} catch {
	// Audit logging failures are non-fatal ← SILENT DATA LOSS
}

\end{lstlisting}



When the disk is full, \texttt{appendFile} throws \texttt{ENOSPC}. The catch block silently discards this error. The audit data is lost with no indication to the operator. Meanwhile, the file system is full, which may cause other parts of the application to fail.

\#\#\# No Existing Rotation Infrastructure

The \texttt{audit.jsonl} file is never:
- Rotated by size (e.g., split at 10 MB)
- Rotated by time (e.g., daily)
- Compressed (old rotations could be gzipped)
- Checked for available disk space before writing
- Capped at a maximum number of rotation files

---

\subsection*{Fix Implementation}
\#\#\# Option A: Size-Based Rotation with Retention Limit (Recommended)



\begin{lstlisting}[language=TypeScript]
// audit.ts — rotated audit logger
import { appendFile, mkdir, rename, stat, readdir, rm } from "fs/promises";
import { createGzip } from "zlib";
import { createReadStream, createWriteStream } from "fs";
import { pipeline } from "stream/promises";

export class AuditLogger {
	private logPath: string;
	private sessionId: string;
	private enabled: boolean;
	private maxSizeBytes: number;
	private maxFiles: number;
	private currentSize: number = 0;

	constructor(
		gitagentDir: string,
		sessionId: string,
		enabled: boolean,
		maxSizeMB: number = 10,   // Default: rotate at 10 MB
		maxFiles: number = 5,      // Default: keep 5 rotated files
	) {
		this.logPath = join(gitagentDir, "audit.jsonl");
		this.sessionId = sessionId;
		this.enabled = enabled;
		this.maxSizeBytes = maxSizeMB * 1024 * 1024;
		this.maxFiles = maxFiles;
	}

	async log(event: string, data: Partial<AuditEntry> = {}): Promise<void> {
		if (!this.enabled) return;

		const entry: AuditEntry = {
			timestamp: new Date().toISOString(),
			session_id: this.sessionId,
			event,
			...data,
		};

		const line = JSON.stringify(entry) + "\n";
		const lineBytes = Buffer.byteLength(line, "utf-8");

		try {
			await mkdir(dirname(this.logPath), { recursive: true });

			// Check size before writing
			if (this.currentSize + lineBytes > this.maxSizeBytes) {
				await this.rotate();
				this.currentSize = 0;
			}

			await appendFile(this.logPath, line, "utf-8");
			this.currentSize += lineBytes;
		} catch (err: any) {
			// Log the failure but don't crash the agent
			console.error(`[audit] Failed to write audit entry: ${err.message}`);
		}
	}

	private async rotate(): Promise<void> {
		// Close current file by renaming with timestamp
		const now = Date.now();
		const rotatedPath = this.logPath + `.${now}`;

		try {
			await rename(this.logPath, rotatedPath);
		} catch {
			return; // File might not exist yet
		}

		// Compress in background (fire-and-forget)
		this.compressAndCleanup(rotatedPath).catch(() => {});
	}

	private async compressAndCleanup(rotatedPath: string): Promise<void> {
		const gzPath = rotatedPath + ".gz";
		try {
			await pipeline(
				createReadStream(rotatedPath),
				createGzip(),
				createWriteStream(gzPath),
			);
			await rm(rotatedPath);
		} catch {
			// Keep uncompressed if compression fails
		}

		// Enforce retention limit
		try {
			const dir = dirname(this.logPath);
			const base = this.logPath.split("/").pop() || "audit.jsonl";
			const files = (await readdir(dir))
				.filter((f) => f.startsWith(base + ".") && f.endsWith(".gz"))
				.sort();

			while (files.length > this.maxFiles) {
				const old = files.shift()!;
				await rm(join(dir, old)).catch(() => {});
			}
		} catch { /* best effort */ }
	}

	// ... (other methods unchanged)
}

\end{lstlisting}



\textbf{What this fixes:}
1. Automatic rotation at configurable size (default 10 MB)
2. Background gzip compression of rotated files
3. Retention limit (default 5 compressed files = \textasciitilde{}50 MB max total)
4. Non-fatal error handling that still logs the failure

\#\#\# Option B: Simple Truncation (Ring Buffer)



\begin{lstlisting}[language=TypeScript]
// audit.ts — ring buffer file
export class AuditLogger {
	private maxEntries: number;
	private entryCount: number = 0;

	async log(...) {
		this.entryCount++;

		// Every N entries, trim the file
		if (this.entryCount % 1000 === 0) {
			await this.trim();
		}
		// ... append as before
	}

	private async trim(): Promise<void> {
		try {
			const content = await readFile(this.logPath, "utf-8");
			const lines = content.split("\n");
			if (lines.length > this.maxEntries) {
				// Keep only the last maxEntries lines
				const trimmed = lines.slice(lines.length - this.maxEntries).join("\n");
				await writeFile(this.logPath, trimmed, "utf-8");
			}
		} catch { /* file may not exist */ }
	}
}

\end{lstlisting}



\#\#\# Option C: Daily Rotation with Size Cap



\begin{lstlisting}[language=TypeScript]
// audit.ts — daily rotation
export class AuditLogger {
	private currentDate: string = "";

	async log(...) {
		const today = new Date().toISOString().slice(0, 10);
		if (today !== this.currentDate) {
			await this.rotateByDate(today);
			this.currentDate = today;
		}
		// ... append
	}

	private async rotateByDate(date: string): Promise<void> {
		const oldPath = this.logPath;
		const newPath = join(dirname(this.logPath), `audit-${date}.jsonl`);
		try {
			await rename(oldPath, newPath);
			// Compress old daily files
		} catch { /* first write of the day */ }
	}
}

\end{lstlisting}



---

\subsection*{Affected Files}
\begin{itemize}
    \item \srcfile{src/audit.ts} — primary affected file
\end{itemize}

\subsection*{Verification}
The fix is verified through unit tests and integration tests that confirm the corrected behavior:

\begin{lstlisting}[language=TypeScript]
// verification/bug-023-verify.ts
import { describe, it, before, after } from "node:test";
import assert from "node:assert";
import { mkdirSync, writeFileSync, readFileSync, rmSync, statSync, existsSync } from "fs";
import { join } from "path";
import { AuditLogger } from "../src/audit";

const TEST_DIR = "/tmp/bug-023-verify";

describe("BUG-023: Audit log rotation", () => {
	before(() => rmSync(TEST_DIR, { recursive: true, force: true }));
	after(() => rmSync(TEST_DIR, { recursive: true, force: true }));

	it("should rotate file when exceeding max size", async () => {
		// Create logger with tiny max size (1 KB) so we can trigger rotation
		const logger = new (AuditLogger as any)(TEST_DIR, "test-session", true, 0.001, 3);
		// Bypass size tracking to test rotation
		await logger.init();

		// Write enough entries to trigger rotation
		for (let i = 0; i < 100; i++) {
			await logger.logToolUse(`tool_${i}`, { arg: "x".repeat(100) });
		}

		// Check that rotation happened
		const dir = TEST_DIR;
		const files = require("fs").readdirSync(dir);
		const auditFiles = files.filter((f: string) => f.startsWith("audit"));
		console.log(`Audit files after writes: ${auditFiles.length}`);
		console.log(`Files: ${auditFiles.join(", ")}`);

		// Should have rotated files
		assert.ok(auditFiles.length > 1, "Should have at least one rotated file");
	});

	it("should enforce retention limit", async () => {
		const logger = new (AuditLogger as any)(TEST_DIR + "/retention", "test", true, 0.001, 2);
		await logger.init();

		for (let i = 0; i < 200; i++) {
			await logger.log(`event_${i}`, { data: "x".repeat(200) });
		}

		const dir = TEST_DIR + "/retention";
		const files = require("fs").readdirSync(dir);
		const gzFiles = files.filter((f: string) => f.endsWith(".gz"));
		console.log(`Compressed audit files: ${gzFiles.length}`);
		assert.ok(gzFiles.length <= 2, `Should retain at most 2 compressed files, got ${gzFiles.length}`);
	});

	it("should not lose recent entries after rotation", async () => {
		const dir = TEST_DIR + "/data";
		const logger = new (AuditLogger as any)(dir, "test", true, 0.01, 3);
		await logger.init();

		// Write 500 entries
		for (let i = 0; i < 500; i++) {
			await logger.logToolUse("test_tool", { iteration: i });
		}

		// The active file should have the most recent entries
		const activeFile = join(dir, "audit.jsonl");
		assert.ok(existsSync(activeFile), "Active audit file should exist");

		const content = readFileSync(activeFile, "utf-8");
		const lines = content.trim().split("\n").filter(Boolean);
		assert.ok(lines.length > 0, "Active file should have recent entries");
		console.log(`Active file has ${lines.length} entries`);
	});
});

\end{lstlisting}

